\section{Failure Is Not Yet a Judgment of the Person}

An objective obligation and personal culpability answer different questions. The first asks what act was due under the governing moral or legal norm. The second asks how a particular person failed, with what knowledge, freedom, intention, and circumstances. Keeping them distinct neither excuses wrongdoing nor permits judgment to outrun the facts.

The \emph{Catechism} says that those who deliberately fail in the Sunday Eucharistic obligation commit a grave sin (CCC 2181). ``Deliberately'' cannot be discarded. Mortal sin requires grave matter, full knowledge, and deliberate consent together (CCC 1857). Knowledge must include the act's sinful character; consent must be sufficiently complete for the act to be a personal choice (CCC 1859). Unintentional ignorance, external pressure, passions, or psychological disturbance can diminish or remove imputability, while chosen ignorance and malice can increase it (CCC 1860).

The resulting propositions are precise. Deliberate omission of the Sunday obligation without serious reason is objectively grave. Not every external absence is a deliberate omission. Not every genuine failure is committed with full knowledge and complete consent. A general article can identify the first proposition; it cannot infer the interior facts needed to pronounce the last judgment upon a named person. Although an act can be judged gravely disordered in itself, judgment of persons is entrusted to God's justice and mercy (CCC 1861).

Canonical and moral effects must also remain distinct. Canon 1247 imposes an obligation. The article does not infer from every violation a canonical delict, penalty, sacramental incapacity, or public status. Moral gravity alone does not create a penal case. Conversely, the absence of a canonical penalty would not make the underlying obligation optional.

In a concrete case, the order of discernment matters: identify the governing law, the person bound, the act required, the material facts, and any applicable excuse, dispensation, or commutation before judging knowledge, consent, or other conditions of culpability. Beginning with subjective innocence can dissolve the norm before its good is seen. Beginning with the external failure and stopping there can mistake law for omniscience. Personal or doubtful cases belong with a pastor, confessor, competent authority, or qualified canonist who has the complete facts.

This is also why obligation should be taught from its object rather than principally from threatened punishment. Sunday protects the worship owed to God, the Eucharistic assembly, Paschal joy, mercy, and rest. The question ``Must I?'' deserves a truthful answer. It also deserves the prior answer to ``What good is being commanded?'' A precept whose good is never named will be heard as arbitrary even when its authority is real.

Yet understanding the good does not suspend obedience until every motive is emotionally compelling. The will can begin by doing what it knows is due. Habit can steady the act; worship can teach desire; grace can heal resistance. The moral life is not divided into a servile stage of rules and a later stage in which love makes justice obsolete. The movement is from an external word truly heard, through willing practice, toward an inward law that gives the act its mature form.
